% Generated by scripts/locobench_paper_examples.py from f224.
\paragraph{Question.} How do economists adjudicate between the claim that the 1970s stagflation originated in oil supply shocks and the claim that it originated in monetary policy errors, and what evidence settled the dispute? (Economics, Social Sciences.)

\paragraph{Claims.} The plan's claims in assertion order, with their factuality labels; claims joined by a ladder's slots are boxed.

\begin{center}\scriptsize
\setlength{\tabcolsep}{2.5pt}
\begin{tabular}{@{}rp{0.76\linewidth}l@{}}
\toprule
& Claim & Label \\
\midrule
\slotclaim{$a_{1}$} & OPEC imposed an oil embargo on the United States in October 1973. & true \\
\slotclaim{$a_{2}$} & The price of crude oil roughly quadrupled between 1973 and early 1974. & true \\
\slotclaim{$a_{3}$} & Hamilton documented a statistical relationship between oil price increases and postwar US recessions. & true \\
$a_{4}$ & Oil supply shocks account for the entire rise in US inflation during the 1970s. & \textbf{false} \\
$a_{5}$ & US consumer price inflation exceeded 5 percent before the October 1973 oil embargo. & true \\
$a_{6}$ & US consumer price inflation remained below 5 percent until the October 1973 oil embargo. & true \\
$a_{7}$ & OPEC production cuts sharply raised the real price of crude oil in 1973 and 1974. & true \\
$a_{8}$ & US money growth accelerated in the years preceding 1973. & true \\
$a_{9}$ & Japan experienced lower consumer price inflation than the United States in 1974. & \textbf{false} \\
$a_{10}$ & The Federal Reserve's policy interest rate stood below the rate of inflation for much of the 1970s. & true \\
$a_{11}$ & Real policy interest rates in the United States were negative during much of the 1970s. & true \\
$a_{12}$ & Clarida, Gali and Gertler estimated that the pre-1979 Federal Reserve policy rule violated the Taylor principle. & true \\
\slotclaim{$a_{13}$} & Paul Volcker became chairman of the Federal Reserve in August 1979. & true \\
\slotclaim{$a_{14}$} & The Volcker tightening pushed the federal funds rate above 19 percent in 1981. & true \\
\bottomrule
\end{tabular}
\end{center}

\paragraph{Template.} The answer as written once, with its slot markers.

\begin{quote}\small
Adjudicating between the supply-shock and the monetary-policy accounts of the 1970s turns on timing, on cross-country comparison, and on the estimated behaviour of the central bank's reaction function. \slotmarker{S1} \slotmarker{S2} That recession regularity indicates that oil supply shocks account for the entire rise in US inflation during the 1970s.

The chronology, however, disciplines any such attribution. Only two descriptions of the pre-embargo record are available, and exactly one of them can be correct: either US consumer price inflation exceeded 5 percent before the October 1973 oil embargo, or it remained below 5 percent until that embargo. The published price indices place the decade in the first case, and inflation already above 5 percent in advance of the embargo cannot be squared with the view that oil supply shocks produced the whole of the decade's inflation. What can be affirmed is that at least one of two impulses was in operation, and perhaps both were: OPEC production cuts sharply raised the real price of crude oil in 1973 and 1974, and US money growth accelerated in the years preceding 1973.

The monetary side of the ledger rests on further evidence. Japan experienced lower consumer price inflation than the United States in 1974. The Federal Reserve's policy interest rate stood below the rate of inflation for much of the 1970s; in other words, real policy interest rates in the United States were negative during much of the decade, and the Japanese comparison for 1974 is precisely the sort of international evidence that shows those American real rates to have been negative. That sustained shortfall of real rates below zero is in turn evidence for the estimate of Clarida, Gali and Gertler that the pre-1979 Federal Reserve policy rule violated the Taylor principle: a central bank that raises its nominal rate by less than one-for-one with inflation will let real rates fall as prices accelerate.

The dispute was settled less by any single regression than by an institutional experiment at the close of the decade. \slotmarker{S3} The lesson drawn from that episode is that a determined central bank can end an inflation whose proximate trigger lay in commodity markets, and the weight of professional judgment shifted accordingly toward the monetary account.
\end{quote}

\paragraph{Fillers.} Each slot's fillers, as admitted by the committee.

\begin{center}\scriptsize
\setlength{\tabcolsep}{2.5pt}
\begin{tabular}{@{}lp{0.80\linewidth}@{}}
\toprule
Filler & Text \\
\midrule
\multicolumn{2}{@{}l}{\slotmarker{S1} chain link slot, $a_{1}$--$a_{2}$} \\
\quad support & OPEC imposed an oil embargo on the United States in October 1973, and that embargo caused the price of crude oil to roughly quadruple between 1973 and early 1974. \\
\quad no link & OPEC imposed an oil embargo on the United States in October 1973. The price of crude oil, which can be traced through posted contract terms, roughly quadrupled between 1973 and early 1974. \\
\addlinespace
\multicolumn{2}{@{}l}{\slotmarker{S2} chain link slot, $a_{2}$--$a_{3}$} \\
\quad support & A quadrupling of that magnitude confirms the statistical relationship Hamilton documented between oil price increases and postwar US recessions. \\
\quad no link & Hamilton, working through the postwar US record, documented a statistical relationship between oil price increases and recessions. \\
\addlinespace
\multicolumn{2}{@{}l}{\slotmarker{S3} control slot, $a_{13}$--$a_{14}$} \\
\quad wording 1 & Paul Volcker became chairman of the Federal Reserve in August 1979, and his arrival led to a tightening that pushed the federal funds rate above 19 percent in 1981. \\
\quad wording 2 & The federal funds rate was driven above 19 percent in 1981 by the tightening that resulted from Paul Volcker's accession to the chairmanship of the Federal Reserve in August 1979. \\
\quad wording 3 & Because Paul Volcker took charge of the Federal Reserve as chairman in August 1979, a tightening followed that carried the federal funds rate past 19 percent in 1981. \\
\bottomrule
\end{tabular}
\end{center}

\paragraph{Rungs.} The five responses of the family, which filler each carries in the edited slots, and gold $\LCS$ at the verdict and the asserted world.

\begin{center}\scriptsize
\setlength{\tabcolsep}{2.5pt}
\begin{tabular}{@{}l>{\raggedright\arraybackslash}p{0.27\linewidth}>{\raggedright\arraybackslash}p{0.30\linewidth}cc@{}}
\toprule
& Ladder, rung & Fillers & \multicolumn{2}{c}{$\LCS$ at world} \\
& & & verdict & asserted \\
\midrule
t0 & chain $\rho_0$ & S1: no link, S2: no link & 0.337 & 0.067 \\
t1 & chain $\rho_1$, control $\rho_0$ & S1: support, S2: no link & 0.398 & 0.086 \\
t2 & chain $\rho_2$ & S1: support, S2: support & 0.478 & 0.115 \\
t3 & control $\rho_1$ & S3: wording 2 & 0.398 & 0.086 \\
t4 & control $\rho_2$ & S3: wording 3 & 0.398 & 0.086 \\
\bottomrule
\end{tabular}
\end{center}

\paragraph{Response t1 in full} (chain $\rho_1$; edited slots highlighted, 440 words).

\begin{quote}\small
Adjudicating between the supply-shock and the monetary-policy accounts of the 1970s turns on timing, on cross-country comparison, and on the estimated behaviour of the central bank's reaction function. \slotfill{S1}{OPEC imposed an oil embargo on the United States in October 1973, and that embargo caused the price of crude oil to roughly quadruple between 1973 and early 1974.} \slotfill{S2}{Hamilton, working through the postwar US record, documented a statistical relationship between oil price increases and recessions.} That recession regularity indicates that oil supply shocks account for the entire rise in US inflation during the 1970s.

The chronology, however, disciplines any such attribution. Only two descriptions of the pre-embargo record are available, and exactly one of them can be correct: either US consumer price inflation exceeded 5 percent before the October 1973 oil embargo, or it remained below 5 percent until that embargo. The published price indices place the decade in the first case, and inflation already above 5 percent in advance of the embargo cannot be squared with the view that oil supply shocks produced the whole of the decade's inflation. What can be affirmed is that at least one of two impulses was in operation, and perhaps both were: OPEC production cuts sharply raised the real price of crude oil in 1973 and 1974, and US money growth accelerated in the years preceding 1973.

The monetary side of the ledger rests on further evidence. Japan experienced lower consumer price inflation than the United States in 1974. The Federal Reserve's policy interest rate stood below the rate of inflation for much of the 1970s; in other words, real policy interest rates in the United States were negative during much of the decade, and the Japanese comparison for 1974 is precisely the sort of international evidence that shows those American real rates to have been negative. That sustained shortfall of real rates below zero is in turn evidence for the estimate of Clarida, Gali and Gertler that the pre-1979 Federal Reserve policy rule violated the Taylor principle: a central bank that raises its nominal rate by less than one-for-one with inflation will let real rates fall as prices accelerate.

The dispute was settled less by any single regression than by an institutional experiment at the close of the decade. \slotfill{S3}{Paul Volcker became chairman of the Federal Reserve in August 1979, and his arrival led to a tightening that pushed the federal funds rate above 19 percent in 1981.} The lesson drawn from that episode is that a determined central bank can end an inflation whose proximate trigger lay in commodity markets, and the weight of professional judgment shifted accordingly toward the monetary account.
\end{quote}

\paragraph{Gold relations of response t1.} Sense, compiled coupling, strength band, and the relation's sign at the verdict (v.) and asserted (a.) world ($+$ support, $-$ conflict, $\cdot$ no factor); a slot id marks a relation drawn by a filler.

\begin{center}\scriptsize
\setlength{\tabcolsep}{2.2pt}
\begin{tabular}{@{}lllllcc@{}}
\toprule
Pair & & Sense & Coupling & Band & v. & a. \\
\midrule
$a_{1}\to a_{2}$ & S1 & Cause-Effect & entailment & strong & $+$ & $+$ \\
$a_{3}\to a_{4}$ &  & Evidence & entailment & weak & $-$ & $+$ \\
$a_{5}\to a_{4}$ &  & Refutation & contradiction & strong & $+$ & $-$ \\
$a_{5}\to a_{6}$ &  & Alternative & exclusive & strong & $-$ & $-$ \\
$a_{7}\to a_{8}$ &  & Disjunction & co-necessity & strong & $+$ & $+$ \\
$a_{9}\to a_{11}$ &  & Evidence & entailment & weak & $+$ & $+$ \\
$a_{10}\to a_{11}$ &  & Restatement & equivalence & strong & $+$ & $+$ \\
$a_{11}\to a_{12}$ &  & Evidence & entailment & moderate & $+$ & $+$ \\
$a_{13}\to a_{14}$ & S3 & Cause-Effect & entailment & strong & $+$ & $+$ \\
\bottomrule
\end{tabular}
\end{center}
